\snsection{JUSTIFICACIÓN}
\justifying

El desarrollo del proyecto responde a la necesidad de mejorar los tiempos de aprendizaje y enseñanza de electrónica de potencia. Esto fue consultado con el titular de la cátedra y remarcó la falta de tiempo para desarrollar todo el contenido académico. Por ello, surge la idea de un sistema didáctico donde se pueda experimentar gran parte del contenido académico de forma práctica y concreta. Además, al ser actualizable mejora aún más las posibilidades de aprendizaje, debido a los constantes avances de las tecnologías, permitiendo agregar nuevos módulos o adaptar los ya existentes sin necesidad de mutar de dispositivo.

En adición, el dispositivo fue desarrollado íntegramente con componentes convencionales y sencillos de conseguir, con el objetivo de reducir costos y mejorar la posibilidad de mantenimiento en caso de ser necesario. Además, contará de un manual de usuario que servirá como guía para utilizar correctamente los módulos y también con preguntas de inquietud para favorecer el incentivo de estudio sobre la temática.


\begin{comment}
Escriba aquí los criterios técnicos, sociales, económicos y razones fundamentales, que acrediten la necesidad de implementar este proyecto.
\end{comment}